\section{Discussion}
\label{sec:discussion}

Our findings demonstrate that vaccination delays function as a structural behavioral barrier, with a perceived disutility comparable to substantial reductions in clinical efficacy. This reinforces earlier intertemporal health investment models which suggest that ``time-to-protection'' is as critical a determinant of health-seeking behavior as a vaccine's safety or performance profile \citep{Grossman1972, Determann2014_PLoS}. By quantifying these trade-offs, this study extends the conceptual framework of hyperbolic discounting to the context of preventive immunization, illustrating that humans assign a disproportionate penalty to even immediate delays in biological protection \citep{Laibson1997, O-Donoghue1999}.

A key contribution of this research is the identification of the moderating roles of institutional trust and biological vulnerability. The observation that low-trust individuals exhibit significantly higher sensitivity to service delays suggesting that institutional confidence serves as a ``behavioral stabilizer'' that mitigates the perceived cost of uncertainty. This aligns with recent global evidence showing that institutional trust is a primary predictor of vaccine acceptance during pandemic volatility \citep{Larson2014, Lazarus2021}. In low-trust environments, administrative friction may exacerbate suspicion and disengagement, potentially creating a negative feedback loop that depresses uptake. Furthermore, the heightened delay burden perceived by medically vulnerable and socioeconomically disadvantaged groups---including older adults and those with chronic conditions---indicates that administrative delays function as a \textit{regressive barrier}, disproportionately penalizing those at the highest biological and socioeconomic risk \citep{Shankar2024, Milkman2021}.

These results have several practical implications for public health implementation. First, our behaviorally-augmented SEIR simulations reveal that even modest behavioral delays can lead to a ``preventive gap,'' significantly expanding the scale of an outbreak. Consequently, increasing operational transparency regarding vaccine availability and wait times is essential to reduce the uncertainty that inflates the perceived burden of delay. Second, our findings support the development of tailored delivery models, such as ``fast-track'' protocols for high-risk populations. By prioritizing access for those who perceive the costs of waiting most acutely (i.e., high-$\kappa$ subgroups), health authorities can improve both the efficiency and the equity of preventive programs \citep{Lancsar2008}. Finally, rebuilding institutional trust through consistent and transparent communication is a prerequisite for the success of long-term immunization strategies \citep{Determann2013_EuroSurveill}.

This study has certain limitations. First, as a discrete choice experiment, the results are based on stated preferences in a hypothetical scenario. While such methods are validated for eliciting latent trade-offs, future research using revealed behavior from registry data is needed to confirm the impact of wait-time interventions on actual coverage. Second, the study was conducted in the specific context of Wuhan; while the structural parameters ($\kappa, \delta$) provide a robust framework, the generalizability of these ``time-trust'' dynamics to health systems with different institutional structures remains to be explored.

In conclusion, behavioral delay is a modifiable determinant of epidemic risk. By addressing administrative friction and institutional trust, public health officials can design more responsive delivery systems. Prioritizing timely and equitable access for vulnerable populations is essential to maximize the public health impact of both existing and future immunization programs.
